% Editable stack of causal dilated convolutions (kernel size 2, dilations
% 1, 2, 4), highlighting the receptive field of the last output.
% Written by Claude (2026-10-14).
\begingroup
\begin{tikzpicture}[x=1cm,y=1cm,
        unit/.style={circle,minimum size=0.42cm,inner sep=0pt,line width=0.8pt},
        idle/.style={unit,draw=drawgray!60,fill=white},
        used/.style={unit,draw=drawblue,fill=drawblue_l},
        tap/.style={drawblue,line width=1.1pt}]
    \def\dx{0.95}
    \def\dy{0.78}
    % every unit, idle by default
    \foreach \l in {0,...,3}{
        \foreach \t in {0,...,7}{
            \node[idle] (u\l-\t) at (\dx*\t,\dy*\l) {};
        }
    }
    % the computation tree of the last output: layer l reads t and t - 2^(l-1)
    \foreach \l/\d/\ts in {3/4/{7}, 2/2/{7,3}, 1/1/{7,5,3,1}}{
        \pgfmathtruncatemacro{\below}{\l-1}
        \foreach \t in \ts{
            \pgfmathtruncatemacro{\back}{\t-\d}
            \draw[tap] (u\below-\t) -- (u\l-\t);
            \draw[tap] (u\below-\back) -- (u\l-\t);
        }
    }
    \foreach \l/\ts in {3/{7}, 2/{7,3}, 1/{7,5,3,1}, 0/{0,...,7}}{
        \foreach \t in \ts{ \node[used] at (u\l-\t) {}; }
    }
    % labels
    \node[anchor=east,font=\small] at (-0.35,0) {Input};
    \node[anchor=east,font=\small] at (-0.35,3*\dy) {Output};
    \foreach \l/\d in {1/1,2/2,3/4}{
        \node[anchor=west,font=\small,text=drawgray] at (7*\dx+0.4,\dy*\l) {$d=\d$};
    }
\end{tikzpicture}
\endgroup
